\chapter{CRONOGRAMA}

% Apresentar crograma de forma visual: tabela ou figura. 

\chapter{CRONOGRAMA}

\begin{table}[H]
\centering
\caption{Cronograma de desenvolvimento do trabalho}
\label{tab:cronograma}
\resizebox{\textwidth}{!}{
\begin{tabular}{|l|c|c|c|c|c|c|}
\hline
\textbf{Atividade} & \textbf{Jun} & \textbf{Jul} & \textbf{Ago} & \textbf{Set} & \textbf{Out} & \textbf{Nov} \\
\hline
Levantamento bibliográfico e fundamentação teórica & X & X &  &  &  &  \\
\hline
Estudo das estruturas de \glspl{bvh} & X & X &  &  &  &  \\
\hline
Definição da arquitetura da aplicação & X & X &  &  &  &  \\
\hline
Implementação do pipeline híbrido de renderização &  & X & X &  &  &  \\
\hline
Implementação do Whitted \textit{ray casting} &  & X & X &  &  &  \\
\hline
Implementação de \textit{soft shadows} e \textit{Importance Sampling} GGX &  &  & X & X &  &  \\
\hline
Implementação das diferentes \glspl{bvh} &  & X & X & X &  &  \\
\hline
Integração das ferramentas de coleta de métricas &  &  & X & X &  &  \\
\hline
Construção das cenas de teste &  &  & X & X &  &  \\
\hline
Execução dos testes e coleta de dados &  &  &  & X & X &  \\
\hline
Análise estatística e interpretação dos resultados &  &  &  &  & X &  \\
\hline
Escrita dos capítulos finais &  &  &  & X & X & X \\
\hline
Revisão e correções do trabalho &  &  &  &  & X & X \\
\hline
Preparação da apresentação final &  &  &  &  &  & X \\
\hline
Entrega final do trabalho &  &  &  &  &  & X \\
\hline
\end{tabular}
}
\end{table}